Die relative Luftfeuchtigkeit ist das Verhältnis der absoluten Luftfeuchtigkeit $\rho_\text{w}$ zur Dichte des gesättigten Wasserdampfs $\rho_\text{s}$ bei derselben Temperatur. Bei \thWO\ ist $\rho_\text{s} = \rsWO$ (FoTaBe S.~192):
\begin{align}
 \varphi &= \fiWF = \frac{\rWk}{\rsW} = \fiW \approx \result{\fiWP{0}{2}}
\end{align}
Kühlt sich die Luft ab, bleibt die absolute Luftfeuchtigkeit gleich. Der Wasserdampf beginnt zu kondensieren, sobald die Sättigungsdichte auf \rWkO\ gesunken ist (Taupunkt). Nach der Tabelle liegt dieser Wert zwischen $\rho_1 = \rsAO$ bei $\vartheta_1 = \thAO$ und $\rho_2 = \rsBO$ bei $\vartheta_2 = \thBO$:
\begin{align}
 \vartheta_\text{T} &= \thTF \\
   &= \thA + \left(\thB-\thA\right)\cdot\frac{\rWk-\rsA}{\rsB-\rsA} \\
   &= \thT \approx \result{\thTP{0}{2}}
\end{align}
